\backhead{FURTHER EXERCISES}{khred}
\begin{exlist}
\item A nozzle of 4 mm ($C_d=0.95$) works at 350 kPa. Find the discharge and the application rate on a $12\times12$ m spacing.\exans{1.14 m$^3$\,h$^{-1}$; 7.9 mm\,h$^{-1}$}
\item Catch cans collect 9, 11, 10, 12, 8, 10, 11, 9, 13, 7 mm. Find Christiansen's CU and the low-quarter distribution uniformity (mean of the lowest quarter over the mean).\exans{86 per cent; 0.78}
\item Wind drift and evaporation lose 12 per cent of the water and the distribution efficiency is 0.85. Find the application efficiency and the gross depth for a net 30 mm.\exans{0.75; 40 mm}
\item A soil has a basic intake of 8 mm\,h$^{-1}$. What is the largest nozzle discharge on an $18\times18$ m spacing that avoids runoff?\exans{2.59 m$^3$\,h$^{-1}$}
\item Ideal jet range: find $v_0^2\sin2\vartheta/g$ for $v_0=25$ m\,s$^{-1}$ at $25^\circ$, and compare with a measured wetted radius of 18 m.\exans{48.8 m; drag and break-up cut it to less than half}
\end{exlist}
\backhead{ESSAY QUESTIONS}{khred}
{\small\itshape Answer each question in about one page, with equations, a sketch and at least one number from this chapter.}\par\vspace{3pt}
\begin{exlist}
\item Explain the overlap of sprinkler patterns and why the spacing must change with the wind.
\item Discuss the losses of sprinkler irrigation (drift, evaporation, runoff) and how each is reduced.
\item Explain the maximum allowable application rate and why it depends on the soil, the slope and the cover.
\item Compare set, travelling and centre-pivot sprinklers for a farm of your choice.
\item Explain the ten textbook checks of the sprinkler layer and what each guards against.
\item Why can a sprinkler with high uniformity still cause runoff?
\end{exlist}
\dykband{EXERCISES IN SWAFL}
{\small Open SWAFL (desktop or online) with the benchmark project. The chapter of the \href{https://github.com/halilunaziru73-creator/Coupled-Water-Fluxes-Pore-Scale-Dynamics-Precision-Irrigation-Mulching-Feedback-Loops/raw/main/Software/v1.1.0/SWAFL_Desktop/docs/manual_v1.1.0/SWAFL_Manual_v1.1.0.pdf}{SWAFL manual} that explains each task is given in brackets; record your values in a table and compare them with what is expected (\emph{Check}).}\par\vspace{3pt}
\begin{exlist}
\item {\color{khblue}[Manual, Chapter 13]} Open \emph{Sprinkler design and pipes}; for one nozzle record the discharge, the wetted radius and the application rate at three pressures.\par\noindent{\small\itshape Check: discharge rises as the square root of pressure.}
\item {\color{khblue}[Manual, Chapter 13]} Change the wind speed and record the uniformity and the drift loss.\par\noindent{\small\itshape Check: uniformity falls and drift rises with wind.}
\item {\color{khblue}[Manual, Chapter 13]} Use \emph{To reduce runoff} on the silt loam and record the change the program proposes.\par\noindent{\small\itshape Check: a lower application rate.}
\item {\color{khblue}[Manual, Chapter 13]} Read the ten textbook checks and verify one of them by hand.\par\noindent{\small\itshape Check: the hand value agrees within 2 per cent.}
\end{exlist}
